\section{Микропрограммное обеспечение микроконтроллера (\texttt{course\_work})}

\subsection{Платформа и зависимости}

Файл конфигурации окружения: \texttt{course\_work/platformio.ini}.

\begin{itemize}
    \item Плата: \texttt{nodemcuv2}.
    \item Фреймворк: \texttt{arduino}.
    \item Основные библиотеки:
    \begin{itemize}
        \item \texttt{Adafruit SSD1306}, \texttt{Adafruit GFX} --- рендеринг пользовательской информации на OLED-дисплее;
        \item \texttt{DHT sensor library} --- взаимодействие с датчиком температуры и влажности DHT11;
        \item \texttt{MQUnifiedsensor} --- считывание и калибровка газового датчика MQ-2;
        \item \texttt{NTPClient} --- синхронизация сетевого времени;
        \item \texttt{FirebaseClient} (mobizt) --- выполнение асинхронных операций чтения/записи в RTDB.
    \end{itemize}
\end{itemize}

Выбранный программный стек ориентирован на обеспечение устойчивой работы в условиях ограниченных аппаратных ресурсов платформы ESP8266: каждая библиотека устраняет конкретный технологический риск в области работы с датчиками, временной синхронизацией, сетевым взаимодействием, визуализацией и облачной синхронизацией.

\subsection{Модульная структура прошивки}

\begin{figure}[htbp]
\centering
\begin{adjustbox}{max width=0.96\textwidth}
\begin{tikzpicture}[node distance=11mm and 11mm]
    \node[depbox] (main) {\texttt{main.cpp}};
    \node[depbox, below left=of main] (app) {\texttt{app\_config.h}\\(тайминги)};
    \node[depbox, below=of main] (sec) {\texttt{secrets.h}\\(Wi-Fi/Firebase)};
    \node[depbox, below right=of main] (sns) {\texttt{sensor\_service}\\\texttt{.h/.cpp}};
    \node[depbox, below=of sec] (netw) {\texttt{network\_service}\\\texttt{.h/.cpp}};
    \node[depbox, below left=of netw] (fb) {\texttt{firebase\_service}\\\texttt{.h/.cpp}};
    \node[depbox, below right=of netw] (disp) {\texttt{display\_service}\\\texttt{.h/.cpp}};

    \draw[deparrow] (main) -- (app);
    \draw[deparrow] (main) -- (sec);
    \draw[deparrow] (main) -- (sns);
    \draw[deparrow] (main) -- (netw);
    \draw[deparrow] (main) -- (fb);
    \draw[deparrow] (main) -- (disp);
    \draw[deparrow] (sns) -- (disp);
    \draw[deparrow] (netw) -- (disp);
\end{tikzpicture}

\end{adjustbox}
\caption{Диаграмма зависимостей модулей прошивки}
\end{figure}

\subsection{Роли модулей}

\begin{itemize}
    \item \texttt{main.cpp} --- точка входа, оркестрация запуска, циклический диспетчер, общая координация сервисов.
    \item \texttt{app\_config.h} --- централизованный набор констант таймингов и пороговых значений.
    \item \texttt{secrets.h} --- параметры доступа к Wi-Fi и Firebase.
    \item \texttt{sensor\_service} --- считывание и калибровка датчиков, нормализация физических измерений.
    \item \texttt{network\_service} --- подключение к сети, поддержание временной синхронизации и проверка корректности временных меток.
    \item \texttt{firebase\_service} --- публикация данных в \texttt{/current}, запись в \texttt{/history}, обработка сигналов активности пользователя.
    \item \texttt{display\_service} --- локальная визуализация текущего состояния и телеметрии.
\end{itemize}

\subsection{Константы и интервалы (\texttt{app\_config})}

\begin{itemize}
    \item \texttt{WIFI\_TIMEOUT\_MS = 15000} --- ограничение времени установления сетевого соединения.
    \item \texttt{SENSOR\_UPDATE\_INTERVAL\_MS = 1000} --- период обновления измерений и экрана.
    \item \texttt{ACTIVITY\_INTERVAL\_MS = 10000} --- период проверки активности пользователя.
    \item \texttt{MIN\_VALID\_EPOCH = 1700000000} --- минимально допустимое значение временной метки после синхронизации.
\end{itemize}

Применяемая схема интервалов формирует иерархию частот: быстрый контур обеспечивает выполнение измерений и взаимодействие с пользователем, медленный контур отвечает за облачную координацию и оптимизацию передачи данных.

\subsection{Инициализация (\texttt{setup()})}

Порядок выполнения:

\begin{enumerate}
    \item \texttt{Serial.begin(115200)}.
    \item \texttt{beginDisplay()}:
    \begin{itemize}
        \item в случае неуспешного выполнения --- переход в бесконечный цикл (жёсткое аварийное завершение).
    \end{itemize}
    \item \texttt{connectWiFiAndTime(SSID, PASSWORD, WIFI\_TIMEOUT\_MS)}:
    \begin{itemize}
        \item установление подключения к сети Wi-Fi;
        \item запуск NTP-клиента (\texttt{pool.ntp.org}, смещение \texttt{0}, UTC);
        \item первичное обновление времени.
    \end{itemize}
    \item При наличии подключения к Wi-Fi: \texttt{beginFirebase(FIREBASE\_HOST)}.
    \item \texttt{beginSensors()}:
    \begin{itemize}
        \item \texttt{dht.begin()};
        \item калибровка датчика MQ-2 (10 итераций, усреднение значения \texttt{R0}).
    \end{itemize}
    \item На OLED-дисплей выводится сообщение \texttt{Init complete / Monitoring...}.
\end{enumerate}

\subsection{Обоснование выбранного порядка инициализации}

\begin{itemize}
    \item \textbf{Display-first:} при возникновении неисправности на уровне индикации устройство немедленно фиксирует отказ и не переходит в «слепой» режим работы.
    \item \textbf{Network-before-cloud:} инициализация Firebase целесообразна исключительно при наличии активного сетевого соединения и корректной временной базы.
    \item \textbf{Sensors-after-core-services:} к моменту начала измерений обеспечена готовность подсистем журналирования, отображения и сетевого взаимодействия.
\end{itemize}

\subsection{Основной цикл (\texttt{loop()})}

\subsubsection*{Непрерывная часть}

\begin{itemize}
    \item \texttt{loopFirebase()} вызывается в каждой итерации цикла для обработки асинхронных ответов и завершённых облачных операций без блокировки основного потока выполнения.
\end{itemize}

\subsubsection*{Периодическая часть (каждые 1 секунду)}

\begin{enumerate}
    \item \texttt{updateNetworkTime()}.
    \item Считывание данных датчиков: \texttt{readSensors()}:
    \begin{itemize}
        \item \texttt{temperature = dht.readTemperature()};
        \item \texttt{humidity = dht.readHumidity()};
        \item \texttt{gas = mq2.readSensor()}.
    \end{itemize}
    \item Получение временных данных:
    \begin{itemize}
        \item \texttt{epoch = getNetworkEpoch()};
        \item \texttt{formattedTime = getNetworkFormattedTime()};
        \item \texttt{timeValid = isNetworkTimeValid(epoch, MIN\_VALID\_EPOCH)}.
    \end{itemize}
    \item Обновление OLED-дисплея: \texttt{renderTelemetry(...)}.
    \item При выполнении условий (Wi-Fi + Firebase ready + valid time):
    \begin{enumerate}
        \item с периодом 10 секунд --- \texttt{checkActivityAndSyncCurrent(...)} и считывание \texttt{/status/last\_viewed};
        \item с периодом 1 час (в первые 60 секунд часового интервала) --- запись истории в \texttt{/history/\{hourStartEpoch\}}.
    \end{enumerate}
\end{enumerate}

\subsection{Логика управления активностью пользователя (\texttt{/current})}

Основной принцип функционирования: \textbf{ветка \texttt{/current} публикуется исключительно при подтверждённой пользовательской активности}.

\begin{enumerate}
    \item Микропрограммным обеспечением сохраняются pending-значения телеметрии и выполняется запрос \texttt{get("/status/last\_viewed")}.
    \item После завершения запроса в \texttt{loopFirebase()} вычисляется значение \texttt{diff = |pendingEpoch - last\_viewed|}.
    \item Если \texttt{diff < 120 секунд}:
    \begin{itemize}
        \item выполняется запись \texttt{/current = \{"t":..., "h":..., "g":...\}}.
    \end{itemize}
    \item В противном случае:
    \begin{itemize}
        \item выполняется удаление узла \texttt{/current}.
    \end{itemize}
\end{enumerate}

Описанный механизм обеспечивает снижение объёма избыточного realtime-трафика и сохраняет предсказуемость поведения системы: «живые» данные публикуются исключительно в период активности клиентского интерфейса.

\subsection{Логика записи истории (\texttt{/history})}

\begin{itemize}
    \item Частота записи: \textbf{1 запись в час}.
    \item Защита от дублирования: \texttt{lastHistoryHourSent}.
    \item Формат пути: \texttt{/history/<epoch начала часа UTC>}.
    \item Полезная нагрузка: \texttt{t, h, g}.
\end{itemize}

Почасовая дискретизация обеспечивает компромисс между детализацией и объёмом хранилища, достаточный для практического анализа динамики на уровне часов и суток.

\subsection{Защита от некорректных показаний датчиков}

Перед передачей данных выполняются следующие проверки:

\begin{itemize}
    \item значения \texttt{temperature} и \texttt{humidity} проходят обработку функцией \texttt{sanitizeReading()}:
    \begin{itemize}
        \item при получении значения \texttt{NaN} производится замена на \texttt{0.0}.
    \end{itemize}
    \item Значение концентрации газа передаётся без изменений (округление до ppm выполняется в клиентском приложении).
\end{itemize}

На OLED-дисплее:

\begin{itemize}
    \item при значении \texttt{NaN} для температуры и влажности отображается \texttt{0}.
\end{itemize}

\subsection{Отображение на OLED-дисплее (\texttt{display\_service})}

\begin{itemize}
    \item Разрешение: \texttt{128x32}.
    \item Штатный режим:
    \begin{itemize}
        \item строка времени (при наличии корректного NTP-времени);
        \item строки \texttt{Tmp}, \texttt{Hum}, \texttt{Gas}.
    \end{itemize}
    \item Режим ожидания синхронизации:
    \begin{itemize}
        \item отображается сообщение \texttt{Time: Syncing...}.
    \end{itemize}
\end{itemize}

\subsection{Вывод по микропрограммному обеспечению}

Микропрограммное обеспечение реализовано в виде неблокирующей сервисной системы с явным разделением быстрых и медленных циклов выполнения. Данное решение обеспечивает устойчивую телеметрию, предсказуемое поведение пользовательского интерфейса и контролируемую нагрузку на облачную инфраструктуру в условиях нестабильного сетевого соединения.
